\section{Developer-Dokumentation}\label{sec:developer_doku}

Die folgende Anleitung beschreibt, wie das Projekt installiert und die verschiedenen Komponenten ausgeführt werden können.
Ziel ist es, das Projekt reproduzierbar und verständlich für Dritte zu machen.

\subsection{Installation}

Das Projekt kann durch Klonen des Repositories und Installation der Dependencies gestartet werden.

\begin{lstlisting}[language=bash]
git clone <repository-url>
cd CounterfactualRegretMinimizationPoker
\end{lstlisting}

Die Dependencies werden mit \texttt{uv sync} installiert:

\begin{lstlisting}[language=bash]
uv sync
\end{lstlisting}

\subsection{Ausführung}

Über die TUI (Text User Interface) können die verschiedenen Komponenten des Projekts gestartet werden.
Die Ausführung unterscheidet sich je nachdem, ob \texttt{mise} installiert ist oder nicht.

Mit \texttt{mise}:
\begin{lstlisting}[language=bash]
mise run tui
\end{lstlisting}

Ohne \texttt{mise}:
\begin{lstlisting}[language=bash]
uv run python src/tui/app.py
\end{lstlisting}

\subsection{Human vs Human Modus}

Im Human vs Human Modus können zwei Spieler über einen Server gegeneinander spielen.
Nach dem Start der TUI wird der Multiplayer-Tab geöffnet, welcher die Konfiguration von Server und Client ermöglicht.

\begin{figure}[h]
\centering
\includegraphics[width=0.9\textwidth]{tui_multiplayer.png}
\caption{TUI Multiplayer-Interface mit Server- und Client-Konfiguration}
\label{fig:tui_multiplayer}
\end{figure}

Die TUI-Oberfläche ist in zwei Bereiche unterteilt: Server und Client.
Im Server-Bereich können das Spiel (Dropdown-Menü), der Port (Eingabefeld) und der Host (Dropdown-Menü) konfiguriert werden.
Mit den Buttons "Start Server" und "Stop Server" wird der Server gestartet beziehungsweise gestoppt.
Der Status-Bereich zeigt an, ob der Server läuft, die lokale IP-Adresse für Remote-Verbindungen und die Connection URL.
Im Client-Bereich ermöglicht der Button "Start Local Client" das Starten eines lokalen Clients.
Für Remote-Verbindungen können die Server-IP, der Port und der Spielername in den entsprechenden Eingabefeldern eingegeben werden.
Der Button "Start Remote Client" startet dann den Remote-Client mit den konfigurierten Parametern.

\begin{figure}[h]
\centering
\includegraphics[width=0.6\textwidth]{connection_info.png}
\caption{Connection Info mit Server-Konfiguration und SSL/WSS-Status}
\label{fig:connection_info}
\end{figure}

Nach dem Start von Server und Client öffnet sich für jeden Spieler ein GUI-Fenster.
Das GUI zeigt die gleiche Poker-Tisch-Oberfläche wie im Agent vs Human Modus, jedoch mit zwei menschlichen Spielern.
Das Layout ist identisch: Board Cards in der Mitte, Pot-Anzeige, Spielerkarten oben und unten, Aktions-Buttons, History-View und Steuerungs-Buttons.
Jeder Spieler sieht das Spiel aus seiner eigenen Perspektive mit seinen eigenen Karten und Aktions-Buttons.

\begin{figure}[h]
\centering
\includegraphics[width=0.9\textwidth]{multiplayer_perspective1.png}
\caption{Multiplayer GUI - Perspektive Spieler 1 (verlierende Hand)}
\label{fig:multiplayer_p1}
\end{figure}

\begin{figure}[h]
\centering
\includegraphics[width=0.9\textwidth]{multiplayer_perspective2.png}
\caption{Multiplayer GUI - Perspektive Spieler 2 (gewinnende Hand)}
\label{fig:multiplayer_p2}
\end{figure}

\subsection{Agent vs Human Modus}

Im Agent vs Human Modus spielt ein menschlicher Spieler gegen einen trainierten Agenten.
Bevor ein Spiel gegen den Agenten gestartet werden kann, muss zunächst eine Strategie für das gewählte Spiel trainiert werden.
Für das Training kann ein beliebiger Algorithmus verwendet werden.
Als Empfehlung gilt der Algorithmus \texttt{discounted\_cfr\_with\_flat\_tree}, da dieser die leistungsstärksten Ergebnisse liefert.

Für die Auswahl des Spiels gibt es zwei Empfehlungen:
Bei wenig verfügbarer Zeit sollte \texttt{twelve\_card\_poker} verwendet werden, da dieses Spiel schneller trainiert werden kann.
Bei mehr verfügbarer Zeit empfiehlt sich \texttt{small\_island\_holdem}, welches komplexer ist und längere Trainingszeiten erfordert.
Beide Spiele wurden im Rahmen dieser Arbeit entwickelt.
Die detaillierten Regeln können im Anhang der Bachelorarbeit im Ordner \texttt{latex/thesis} nachgeschlagen werden.

Das Training wird über die TUI im Tab "Training Queue" konfiguriert.
Dort kann ein neuer Training Task erstellt werden, bei dem das Spiel, der Algorithmus, die Anzahl der Iterationen und weitere Parameter festgelegt werden können.
Das Feld "Early Stop Exploitability" ermöglicht es, das Training automatisch zu beenden, sobald ein bestimmter Exploitability-Wert erreicht wird.
Dies kann Zeit sparen, wenn bereits eine ausreichend gute Strategie erreicht wurde.

\begin{figure}[h]
\centering
\includegraphics[width=0.8\textwidth]{training_task.png}
\caption{TUI-Dialog zur Erstellung eines neuen Training Tasks}
\label{fig:training_task}
\end{figure}

Nach dem Start des Trainings kann der Trainingsverlauf im Tab "Current Task" überwacht werden.
Dieses Fenster zeigt wichtige Informationen wie den aktuellen Fortschritt, die verstrichene Zeit, die geschätzte verbleibende Zeit sowie eine Tabelle mit Best Response Evaluations.
Anhand dieser Metriken kann überprüft werden, ob das Training wie erwartet funktioniert und ob die Exploitability-Werte sinken.

\begin{figure}[h]
\centering
\includegraphics[width=0.9\textwidth]{current_task.png}
\caption{TUI Current Task Fenster zur Überwachung des Trainingsverlaufs}
\label{fig:current_task}
\end{figure}

Nach Abschluss des Trainings kann im Tab "Models" ein trainiertes Modell ausgewählt werden, um gegen den Agenten zu spielen.
Falls die Liste der Modelle nicht aktuell ist, kann sie mit der Taste \texttt{r} aktualisiert werden.
Nach der Auswahl eines Modells durch Anklicken kann mit der Taste \texttt{g} ein Spiel gegen den Agenten gestartet werden.
Daraufhin öffnet sich automatisch das GUI-Fenster.

\begin{figure}[h]
\centering
\includegraphics[width=0.9\textwidth]{models_tab.png}
\caption{TUI Models Tab zur Auswahl eines trainierten Modells}
\label{fig:models_tab}
\end{figure}

Das GUI zeigt eine Poker-Tisch-Oberfläche mit grünem Filz-Hintergrund.
In der Mitte des Tisches werden die Board Cards (Community Cards) angezeigt, die von allen Spielern geteilt werden.
Links neben den Board Cards wird der aktuelle Pot-Wert als Zahl angezeigt, begleitet von visuellen Chip-Stapeln, die den Pot-Betrag darstellen.

Am oberen Rand des Fensters befindet sich der Gegenspieler (im Agent vs Human Modus der "Strategy Agent", im Multiplayer Modus der andere menschliche Spieler).
Die Karten des Gegenspielers werden angezeigt, sobald die Hand beendet ist.
Unter den Karten wird die Hand-Stärke des Gegenspielers angezeigt (z.B. "High Card", "Pair", "Two Pair", "Straight").

Am unteren Rand befindet sich der eigene Spieler mit seinen Hole Cards.
Auch hier wird die Hand-Stärke unter den eigenen Karten angezeigt.
Darunter befinden sich die Aktions-Buttons: "Check", "Bet", "Call" (mit dem zu zahlenden Betrag) und "Fold".
Diese Buttons sind nur aktiv, wenn der Spieler am Zug ist.

Im oberen rechten Bereich befinden sich die Steuerungs-Buttons: "New Hand" zum Starten einer neuen Hand und "Beenden" zum Beenden des Spiels.
Darunter befindet sich der History-View, der alle Aktionen während der aktuellen Hand chronologisch auflistet.
Am Ende jeder Hand wird das Ergebnis angezeigt (z.B. "You wins 100!" oder "Opponent wins 100!") sowie die Payoff-Informationen.

\begin{figure}[h]
\centering
\includegraphics[width=0.9\textwidth]{gui_agent_vs_human.png}
\caption{GUI für Agent vs Human Modus mit abgeschlossenem Spiel}
\label{fig:gui_agent_vs_human}
\end{figure}

Das GUI ermöglicht es dem menschlichen Spieler, verschiedene Strategien zu testen, einschließlich Bluffs.
Das History-Log dokumentiert alle Aktionen und zeigt das Ergebnis jeder Hand an.

\begin{figure}[h]
\centering
\includegraphics[width=0.9\textwidth]{gui_bluff.png}
\caption{GUI mit erfolgreichem Bluff: Der menschliche Spieler gewinnt durch einen Bet, der den Agenten zum Folden bringt}
\label{fig:gui_bluff}
\end{figure}
